% !TEX root = ../00_Main/Main.tex
\documentclass[../00_Main/Main.tex]{subfiles}

\begin{document}

Esta tesis compara arquitecturas de aprendizaje por refuerzo (RL,
\emph{Reinforcement Learning}) en una simulaci\'on de pandemia donde el
agente estudiado distribuye un presupuesto fijo entre ciudades que se presentan sucesivamente, sin conocer cu\'antas
restan ni cu\'al ser\'a su severidad. El entorno, \textbf{mPES}, extiende el \emph{Pandemic
Experiment Scenario} (PES) del BCI-NE Lab de la Universidad de
Essex~\citep{BCINE2022}.

\subsection{Motivaci\'on}\label{sec:motivacion}

Durante la pandemia de COVID--19, gobiernos y hospitales debieron
asignar recursos escasos en condiciones en que cada asignaci\'on
modificaba la evoluci\'on posterior del sistema sanitario. En esta tesis
la pandemia no es el objeto de estudio, sino un escenario controlado que
permite medir qu\'e modelos toman mejores decisiones y c\'omo se
comportan fuera de las condiciones de entrenamiento.

\subsection{Definici\'on del problema}\label{sec:problema}

En cada paso se presenta una ciudad y el agente le asigna entre $0$ y
$10$ unidades de recurso. La severidad de cada ciudad crece en un factor $\beta$ por paso y
cada unidad asignada la reduce en una cantidad $\alpha$
(Ecuaci\'on~\eqref{eq:transicion}).
Esta regla simplifica los modelos de la epidemiolog\'ia
computacional~\citep{Kuhl2021}, desde los compartimentales hasta
simulaciones espaciales de COVID-19 basadas en agentes, como las de Parisi
\emph{et al.}~\citep{Parisi2021}: conserva s\'olo el
crecimiento de la severidad y el efecto de los recursos.

\subsection{Preguntas e hip\'otesis}\label{sec:hypothesis}

\begin{enumerate}
    \item ?`Qu\'e modelo individual obtiene el mayor desempe\~no, y lo
    conserva cuando cambian la severidad o la longitud de las
    secuencias?
    \item ?`Alg\'un ensamble de modelos ya entrenados supera al mejor
    modelo individual y, en tal caso, con qu\'e regla de combinaci\'on?
\end{enumerate}

Para cada pregunta se plantea una hip\'otesis:

\begin{quote}
\textbf{H1.} \emph{Entre los modelos individuales evaluados, el
Transformer causal alcanza el mayor desempe\~no medio y lo conserva en
los escenarios de generalizaci\'on.}

\textbf{H2.} \emph{Los ensambles de redes neuronales con votaci\'on
basada en la confianza derivada de la entrop\'ia de Shannon superan el
desempe\~no del mejor modelo individual.}
\end{quote}

\subsection{Contribuciones}

\begin{enumerate}
    \item Una comparaci\'on de seis modelos de RL, evaluados con los
    mismos entorno, protocolo y cantidad de secuencias, con
    hiperpar\'ametros de entrenamiento ajustados por optimizaci\'on
    bayesiana.
    \item Seis reglas para combinar esos modelos sin reentrenarlos, en
    las que cada modelo se pondera seg\'un la entrop\'ia de su
    distribuci\'on de acciones.
    \item Una evaluaci\'on en $21$ escenarios que cambian la severidad,
    la longitud de las secuencias o ambas, en un \emph{framework}
    reproducible.
\end{enumerate}

\biblio % Needed for referencing to working when compiling individual subfiles - Do not remove

\end{document}
